
\documentclass{article} % For LaTeX2e
\usepackage{iclr2026_conference,times}
\iclrfinalcopy
% Optional math commands from https://github.com/goodfeli/dlbook_notation.
\input{math_commands.tex}

\usepackage{hyperref}
\usepackage{url}
\usepackage{booktabs} 
\usepackage{graphicx}
\usepackage{multirow}
\usepackage{amsmath}
\usepackage{mathtools}
\usepackage{tcolorbox}
\usepackage{enumitem}
\usepackage{geometry}
\usepackage{caption}
\usepackage{subcaption}
\usepackage{listings}
\tcbuselibrary{skins,listings,breakable,listingsutf8}
\usepackage{wrapfig}
\usepackage{xspace}

\definecolor{darkblue}{rgb}{0, 0, 0.5}
\hypersetup{colorlinks=true, citecolor=darkblue, linkcolor=darkblue, urlcolor=darkblue}

\newcommand{\methodname}{\textsc{Intuitor}\xspace}

\newcommand{\xuandong}[1]{\textcolor{red}{[Xuandong: #1]}}
\newcommand{\dawn}[1]{\textcolor{red}{[Dawn: #1]}}
\newcommand{\todo}[1]{\textcolor{red}{[TODO: #1]}}
\newcommand{\blue}[1]{\textcolor{blue}{#1}}


\title{Learning to Reason without External Rewards}

% Authors must not appear in the submitted version. They should be hidden
% as long as the \iclrfinalcopy macro remains commented out below.
% Non-anonymous submissions will be rejected without review.

\author{%
    % , Xuandong Zhao, Aosong Feng, Sergey Levine, Dawn Song
Xuandong Zhao\thanks{Equal contribution.}\\
UC Berkeley\\
{\small \texttt{xuandongzhao@berkeley.edu}} \\
\And
Zhewei Kang\footnotemark[1] \\
UC Berkeley\\
{\small \texttt{waynekang@berkeley.edu}} \\
\And
Aosong Feng\\
Yale University\\
{\small \texttt{aosong.feng@yale.edu}}\\
\AND 
Sergey Levine\\
UC Berkeley\\
{\small \texttt{svlevine@berkeley.edu}}\\
\And
Dawn Song\\
UC Berkeley\\
{\small \texttt{dawnsong@berkeley.edu}}\\
}


% The \author macro works with any number of authors. There are two commands
% used to separate the names and addresses of multiple authors: \And and \AND.
%
% Using \And between authors leaves it to \LaTeX{} to determine where to break
% the lines. Using \AND forces a linebreak at that point. So, if \LaTeX{}
% puts 3 of 4 authors names on the first line, and the last on the second
% line, try using \AND instead of \And before the third author name.

\newcommand{\fix}{\marginpar{FIX}}
\newcommand{\new}{\marginpar{NEW}}

%\iclrfinalcopy % Uncomment for camera-ready version, but NOT for submission.
\begin{document}


\maketitle

% \title{Learning to Reason without External Rewards}
\begin{abstract}
Training large language models (LLMs) for complex reasoning via Reinforcement Learning with Verifiable Rewards (RLVR) is effective but limited by reliance on costly, domain-specific supervision. We explore Reinforcement Learning from Internal Feedback (RLIF), a framework that enables LLMs to learn from intrinsic signals without external rewards or labeled data. We propose \methodname, an RLIF method that uses a model's own confidence—termed \emph{self-certainty}—as its sole reward signal. \methodname replaces external rewards in Group Relative Policy Optimization (GRPO) with self-certainty scores, enabling fully unsupervised learning. Experiments demonstrate that \methodname matches GRPO's performance on mathematical benchmarks while achieving better generalization to out-of-domain tasks like code generation, without requiring gold solutions or test cases. Our findings show that intrinsic model signals can drive effective learning across domains, offering a scalable alternative to RLVR for autonomous AI systems where verifiable rewards are unavailable. Code is available at \url{https://github.com/sunblaze-ucb/Intuitor}.
\end{abstract}

\begin{figure}[h]
    \centering
    \includegraphics[width=0.92\linewidth]{fig/page1.pdf}
    \vspace{-0.5em}
    \caption{Overview of RLIF and \methodname's Performance. Left: RLIF, a paradigm where LLMs learn from intrinsic signals generated by the model itself, without external supervision. Right: Performance comparison of Qwen2.5-3B Base, GRPO, and \methodname (our RLIF instantiation). Both GRPO and \methodname are trained on the MATH dataset. \methodname achieves comparable performance to GRPO on in-domain mathematical benchmarks (GSM8K, MATH500) and demonstrates better generalization to out-of-domain code generation tasks (LiveCodeBench-v6, CRUXEval). Part of the illustration was generated by GPT-4o.
    }
    \label{fig:main}
\end{figure}


\section{Introduction}

Reinforcement learning has become essential for enhancing large language model capabilities. Early work focused on Reinforcement Learning from Human Feedback (RLHF), which aligns model outputs with human values through reward models trained on preference data~\citep{ouyang2022training}. Recent advances in Reinforcement Learning with Verifiable Rewards (RLVR) replace learned reward models with automatically verifiable signals, such as exact answer matching in mathematical problem-solving, demonstrating improved reasoning capabilities in models like DeepSeek-R1 \citep{guo2025deepseek,lambert2024t}.

Despite these successes, both RLHF and RLVR face fundamental limitations that constrain their broader applicability. RLHF requires extensive human annotation, making it expensive and potentially biased~\citep{gao2023scaling}. RLVR, while avoiding learned reward models, demands domain-specific verifiers and gold-standard solutions. In mathematics, this requires expert annotation of solutions; in code generation, it necessitates comprehensive test suites and execution environments~\citep{liu2023your,code-r1,team2025kimi,xiaomi2025mimo}. These requirements limit RLVR to carefully curated domains and complicate deployment in open-ended scenarios. Moreover, outcome-oriented verifiable rewards limit transferability to other domains. These challenges motivate exploration of more general and scalable reward paradigms, leading to a critical research question:
\begin{center}
\emph{Can LLMs enhance their reasoning abilities by relying solely on intrinsic, self-generated signals, without recourse to external verifiers or domain-specific ground truth?}
\end{center}

In this paper, we introduce and explore such a paradigm: \emph{Reinforcement Learning from Internal Feedback (RLIF)}, where models optimize intrinsic feedback to improve performance without external rewards or supervision. The motivation for RLIF extends to future scenarios where models develop superhuman capabilities that become difficult for humans to evaluate directly \citep{burns2023weak}, requiring self-improvement through intrinsic mechanisms \citep{oudeyer2007intrinsic}.

Under the RLIF paradigm, we propose \methodname, a novel reinforcement learning approach leveraging a model's own confidence as an intrinsic reward. This builds on observations that LLMs exhibit lower confidence on difficult problems~\citep{farquhar2024detecting, kuhn2023semantic, kang2024unfamiliar, kang2025scalable}; optimizing for confidence should improve reasoning capabilities. Specifically, we use self-certainty \citep{kang2025scalable}, the average KL divergence between the model's output distribution and a uniform distribution, as our confidence measure. This metric has proven useful for distinguishing high-quality responses from flawed ones \citep{kang2025scalable, ma2025reasoning}. Building on this insight, \methodname guides learning through self-generated signals, eliminating the need for external supervision or handcrafted rewards. The implementation of \methodname is simple, efficient, and effective: we replace the verifiable reward signal in existing RLVR frameworks, specifically Group Relative Policy Optimization (GRPO) \citep{shao2024deepseekmath}, with self-certainty scores, using the same policy gradient algorithm.

Our experiments demonstrate promising results. On the MATH dataset \citep{hendrycks2021measuring} with Qwen2.5-3B base \citep{qwen2.5}, \methodname matches the performance of GRPO without relying on any gold answers. As \methodname rewards the generation trajectory rather than only the end result, it generalizes more effectively: training a Qwen2.5-3B base model on MATH yields a \(65\%\) relative improvement on LiveCodeBench Code generation task \citep{jain2024livecodebench} versus no improvement for GRPO, and a \(76 \%\) gain on CRUXEval-O \citep{gu2024cruxeval} compared with \(44 \%\) for GRPO. Additionally, when we fine-tune the Qwen2.5-1.5B base model with \methodname on the MATH corpus, a model that originally produces repetitive content and scores \(0\%\) on LiveCodeBench learns to emit coherent reasoning chains and well-structured code, reaching \(9.9\%\) accuracy after the tuning. Beyond the Qwen family, experiments with Llama~\citep{meta2024llama3} and OLMo~\citep{olmo20242olmo2furious} models also show impressive gains, underscoring the strong generalization capabilities of \methodname. As \methodname requires only a clear prompt and no verifiable reward, it is broadly applicable across tasks, providing fresh evidence that pretrained LLMs possess richer latent behavioral priors than previously recognized. 

Our contributions can be summarized as follows:
% \begin{itemize}[leftmargin=*,topsep=0pt,itemsep=0pt]
\begin{itemize}[leftmargin=*, nosep]
  \item We introduce and explore Reinforcement Learning from Internal Feedback (RLIF), a novel reinforcement learning paradigm enabling LLMs to improve reasoning skills by leveraging intrinsic,  self-generated signals, without reliance on external supervision or labeled data.
  \item We introduce \methodname, an RLIF-based method that utilizes a model's own internal confidence measure—termed \emph{self-certainty}—as the sole intrinsic reward. 
  \item We demonstrate that \methodname matches supervised RL performance on in-domain tasks and achieves competitive, sometimes better out-of-domain generalization. We uncover emergent structured reasoning and enhanced instruction-following capabilities induced by intrinsic rewards.
\end{itemize}

% Experiments demonstrate that Intuitor matches GRPO's performance on mathematical benchmarks and achieves competitive, sometimes better, generalization to out-of-domain tasks like code generation, without requiring gold solutions or test cases.

\section{Related Work}

\paragraph{Reinforcement Learning from Human Feedback (RLHF).}  RL has become instrumental in refining LLMs. Early pivotal work centered on Reinforcement Learning from Human Feedback (RLHF) \citep{ouyang2022training}, which aligns LLMs with human values by training a reward model on human preference data. While effective, RLHF is often resource-intensive due to the need for extensive human annotation \citep{touvron2023llama}. Subsequent innovations like Direct Preference Optimization (DPO) \citep{rafailov2023direct} aimed to simplify this by directly training models on preferences. The reliance on human-generated or model-approximated human preferences poses scalability challenges and introduces potential biases from the reward model itself \citep{gao2023scaling}.

\paragraph{Reinforcement Learning with Verifiable Rewards (RLVR).} RLVR emerged as a powerful alternative, particularly for tasks with clear correctness criteria like mathematical reasoning and code generation \citep{hu2025open,team2025kimi,xiaomi2025mimo}. RLVR utilizes rule-based verification functions, such as exact answer matching \citep{guo2025deepseek,team2025kimi,xiaomi2025mimo,jaech2024openai}, to provide reward signals, thereby avoiding the complexities and potential pitfalls of learned reward models. This approach has sparked significant advances, with models like DeepSeek-R1 \citep{guo2025deepseek} achieving state-of-the-art reasoning capabilities. The development of robust policy optimization algorithms like GRPO \citep{shao2024deepseekmath} and its variants \citep{deepcoder2025,liu2025understanding} has further solidified RLVR's success. Nevertheless, RLVR's applicability is largely confined to domains where verifiable gold solutions or exhaustive test cases can be constructed, and its predominant focus on outcome-based rewards can limit generalization to dissimilar tasks or those requiring nuanced, process-oriented feedback. 


\paragraph{Intrinsic Signals and Self-Play in LLM Optimization.} Self-play and intrinsic rewards enable autonomous model improvement. Methods like SPIN \citep{chen2024spin} and Self-Rewarding LMs \citep{yuan2024selfrewarding} use the model itself for feedback. Earlier work like STaR \citep{zelikman2022star} relies on outcome evaluation, while others explore procedural generalization \citep{poesia2024minimo,cheng2024spag}. Concurrent works such as Genius, TTRL, SRT, and Absolute Zero \citep{xu2025genius,zuo2025ttrl,shafayat2025can,zhao2025absolute} leverage unlabeled queries for RL but are often restricted to specific task distributions. \citet{song2025mind} examine LLM self-improvement through the generation–verification gap, while \citet{huang2025self} study it through the lens of sharpening dynamics. \methodname aligns with this direction, offering a lightweight, general-purpose approach using self-certainty as a confidence-based intrinsic reward, enabling single-agent RL across diverse tasks without explicit feedback or gold labels.

\section{Method}

\subsection{Reinforcement Learning from Internal Feedback (RLIF)}
To overcome the limitations of RLHF's costly human annotation and RLVR's domain-specific supervision, we propose Reinforcement Learning from Internal Feedback (RLIF). Instead of depending on external evaluation, RLIF uses the model's own assessment of its outputs as feedback. This offers several advantages: it reduces reliance on supervision infrastructure, provides task-agnostic reward signals, and supports learning in domains where external verification is unavailable. The optimization objective for policy $\pi_\theta$ is:
\begin{equation}\label{eq:intuitor}
\max_{\pi_\theta} \mathbb{E}_{o \sim \pi_\theta(q)} \left[u(q, o) - \beta \mathrm{KL}[\pi_\theta(o|q) \| \pi_{\text{ref}}(o|q)] \right]
\end{equation}
where $q$ is an input query, $o$ is the generated output, $\pi_{\text{ref}}$ is an initial reference policy, and $\beta$ controls the KL divergence to prevent excessive deviation from $\pi_{\text{ref}}$. Here, $u(q,o)$ is an intrinsic signal derived from the model's internal state or computation, rather than external verification. The key challenge lies in identifying intrinsic signals that correlate with output quality and can effectively guide learning.

Concurrent research explores related concepts within the RLIF paradigm. For example, Entropy Minimized Policy Optimization (EMPO) \citep{zhang2025right} minimizes LLM predictive entropy on unlabeled questions in a latent semantic space. SEED-GRPO \citep{chen2025seed} uses the semantic entropy of generated sequences, combined with ground truth rewards, to modulate policy updates. Reinforcement Learning with a Negative Entropy Reward (EM-RL) \citep{agarwal2025unreasonable} employs a reward signal based solely on the negative sum of token-level entropy, akin to REINFORCE but without labels. These methods highlight the growing interest and potential of leveraging intrinsic signals for LLM training under the RLIF framework.

\subsection{\methodname: Policy Optimization with Self-Certainty}

We propose \methodname, a novel RLIF method that utilizes a model's own confidence as the sole intrinsic reward signal $u(q,o)$. Our choice of model confidence as the intrinsic reward is motivated by observations that LLMs often exhibit lower confidence when encountering unfamiliar tasks or lacking sufficient knowledge \citep{kang2024unfamiliar}. Conversely, higher confidence frequently correlates with correctness. By rewarding increased self-confidence, \methodname encourages to iteratively ``practice'' and refine its reasoning pathways until it becomes more confident in its outputs.

We adopt the self-certainty metric from \citet{kang2025scalable}, defined as the average KL divergence between a uniform distribution $U$ over the vocabulary $\mathcal{V}$ and the model's next-token distribution:
\begin{align}
\textbf{Self-certainty}(o|q)  \coloneqq \frac{1}{|o|}\sum_{i=1}^{|o|} \mathrm{KL}(U\parallel p_{\pi_\theta}(\cdot|q,o_{<i}))  = -\frac{1}{|o|\cdot|\mathcal{V}|}\sum_{i=1}^{|o|}\sum_{j=1}^{|\mathcal{V}|} \log\left(|\mathcal{V}|\cdot p_{\pi_\theta}(j|q,o_{<i})\right) \label{eq:self_certainty}
\end{align}
where $o_{<i}$ are the previously generated tokens and $p(j|q,o_{<i})$ is the model's predicted probability for token $j$ at step $i$. Higher self-certainty values indicate greater confidence. 


Self-certainty is related to a KL divergence where the model's prediction is the second argument, $\mathrm{KL}(U\parallel p_{\pi_\theta})$. This contrasts with entropy (or reverse KL divergence from uniform).
Critically, self-certainty is reported to be less prone to biases towards longer generations, a common issue with perplexity or entropy-based measures \citep{fang2024wrong, kang2025scalable}, making it a potentially more reliable indicator of intrinsic confidence. \citet{kang2025scalable} demonstrate that self-certainty is effective for selecting high-quality answers from multiple candidates, and uniquely among different confidence measures, its utility improves with more candidates. Optimizing for self-certainty thus encourages the model to generate responses that it deems more convincing. The RL process can achieve this by, for instance, guiding the model to produce more detailed reasoning steps, thereby increasing the model's conviction in its final answer. This mechanism is more nuanced than simply increasing the probability of the most likely output; it involves modifying the generation process itself to build confidence.

\begin{figure}[t]
    \centering
    \vspace{-1em}
    \includegraphics[width=0.92\linewidth]{fig/spo.pdf}
    % \vspace{-0.5em}
    \caption{\methodname simplifies the training strategy by leveraging self-certainty (the model's own confidence) as an intrinsic reward to incentivize reasoning abilities without external supervision.}
    % \caption{\textbf{\methodname uses self-certainty as an intrinsic reward to optimize reasoning without external supervision.} The figure illustrates the GRPO-based training loop in which, for each query, multiple candidate solutions are sampled, scored using self-certainty computed from the model's token-level distributions, and then used to compute advantages for a KL-regularized policy update. This pipeline replaces external verifiable rewards with internal feedback while keeping the underlying optimization algorithm unchanged.}
    \vspace{-1.em}
\label{fig:spo_overview}
\end{figure}

To optimize the objective in Equation \ref{eq:intuitor}, various policy gradient algorithms can be employed. Informed by the recent success in models such as DeepSeek-R1 \citep{guo2025deepseek} and its widespread adoption of GRPO in open-source projects, we utilize GRPO to optimize for self-certainty. The overall pipeline for this GRPO-based instantiation of \methodname is illustrated in Figure~\ref{fig:spo_overview}.

The core idea behind the optimization is to sample multiple candidate outputs for a given query and use their relative rewards to estimate advantages for policy updates. For each query $q \sim P(Q)$, GRPO samples a group of $G$ outputs ${o_1, \ldots, o_G}$ using a behavior policy $\pi_{\theta_{\text{old}}}$ (e.g., a previous iteration or the SFT model). The target policy $\pi_\theta$ is then optimized by maximizing:
\begin{align*}
& \mathcal{J}_{\text{GRPO}}(\theta) =  \mathbb{E}_{\substack{q \sim P(Q), \\  \{o_i\}_{i=1}^G \sim \pi_{\theta_{\text{old}}}(\cdot|q)}}\left[ 
\frac{1}{G}\sum_{i=1}^G\frac{1}{|o_i|}\sum_{t=1}^{|o_i|}
\left(\min \left[c_{i,t}(\theta) \hat{A}_{i,t},\;\operatorname{clip}_{\varepsilon}\left(c_{i, t}(\theta)\right)\hat{A}_{i,t}\right] -\beta \mathbb{D}_{\mathrm{KL}}\bigl(\pi_\theta\Vert\pi_{\mathrm{ref}}\bigr) \right) \right]
\end{align*}
where $c_{i,t}(\theta) = \frac{\pi_\theta(o_{i,t}\mid q,o_{i,<t})} {\pi_{\theta_{\mathrm{old}}}(o_{i,t}\mid q,o_{i,<t})}$ is the importance weight, $\operatorname{clip}_{\varepsilon}$ is the function that clips to $[1-\varepsilon, 1+\varepsilon]$.
Hyperparameters $\epsilon$ (for clipping) and $\beta$ (for KL penalty strength) control stability and exploration, and $\hat{A}_{i,t}$ is the advantage estimate.

\textbf{Integration of Self-Certainty. }
The key innovation in \methodname is replacing external rewards with self-certainty scores in GRPO's advantage computation. Specifically, each output $o_i$ is scored by:
\begin{equation}
u_i = \text{Self-certainty}(o_i|q), \quad
\hat{A}_{i,t} = \frac{u_i - \text{mean}(\{u_1, u_2, \cdots, u_G\})}
{\text{std}(\{u_1, u_2, \cdots, u_G\})}.
\label{eq:spo_reward}
\end{equation}
This formulation enables the policy to favor outputs that the model itself considers more confident. The complete \methodname training pipeline operates by sampling multiple candidate outputs for each query, computing self-certainty scores for each candidate, using these scores to estimate advantages within the group, and updating the policy to increase the likelihood of generating high-confidence outputs. This process requires no external supervision, making it a self-reinforcing learning loop.

\section{Experimental Setup}

\textbf{Training Setup. }
Both GRPO and \methodname are trained with the Open-R1 framework \citep{openr1} on the training split of the MATH dataset \citep{hendrycks2021measuring}, which contains 7,500 problems. We use Qwen2.5-1.5B and Qwen2.5-3B \citep{qwen2.5} as backbone models, with a chat-based prompting format throughout. Given the models’ initially weak instruction-following abilities, we do not require them to disentangle intermediate reasoning from final answers. Each update processes 128 problems, generating 7 candidate solutions per problem, with a default KL penalty of \(\beta = 0.005\). For a fair comparison, GRPO and \methodname share identical hyperparameters (see Appendix) without additional tuning. We also evaluate a GRPO variant, denoted GRPO-PV in Table~\ref{tab:big}, which uses plurality voting\footnote{Self-consistency uses a plurality rule, selecting the most frequent answer even without majority support, while majority voting requires \(> 50\%\) support and otherwise yields no winner \citep{de2014essai}.} as a proxy for ground truth. This follows the approach from TTRL \citep{zuo2025ttrl}, which shows that self-consistency-based rewards can match the performance of golden answers when training on inference data.

\textbf{\methodname for Code Generation (\methodname-Code). }
To assess generalization beyond mathematical reasoning, we apply \methodname to the Codeforces code generation dataset \citep{li2022competition}. This variant, denoted \methodname-Code in Table~\ref{tab:big}, modifies the setup as follows: the number of sampled completions per problem is increased to 14; the learning rate is reduced from \(3 \times 10^{-5}\) to \(1 \times 10^{-5}\); and the KL penalty is increased to \(\beta = 0.01\). For simplicity, we limit the run to 50 steps, utilizing a total of 3,200 problems.

\textbf{Evaluation. }
Evaluations generally use the same chat-style prompting format as in training, except for MMLU-Pro \citep{wang2024mmlu}, where we follow the benchmark’s original prompt format. Greedy decoding is used for all completions. Experiments were conducted on NVIDIA A100 GPUs, each with 40GB of memory. We evaluate performance on the following benchmarks (1) \emph{Math reasoning}: MATH500 and GSM8K, using the \texttt{lighteval} library \citep{lighteval}.
(2) \emph{Code reasoning}: CRUXEval-O \citep{gu2024cruxeval}, using the \texttt{ZeroEval} framework \citep{Lin_ZeroEval_A_Unified_2024}, and LiveCodeBench v6 (LCB) \citep{jain2024livecodebench}.
(3) \emph{Instruction following}: AlpacaEval 2.0 with length-controlled win rates \citep{dubois2024length}, judged by GPT-4.1 \citep{openai2025gpt4.1}.

\begin{table}[t]
  \centering
  \caption{Performance comparison of various methods on reasoning and instruction-following benchmarks. The \methodname-Code variant is trained on Codeforces data with a smaller learning rate and fewer training steps. All evaluations are obtained with the chat inference template, except for MMLU-Pro.}
  \vspace{-1em}
  \label{tab:big}
  \resizebox{0.98\linewidth}{!}{
  \begin{tabular}{lcccccccc}
    \toprule
    Model           & Training Data & GSM8K  & MATH500 & LCB    & CRUX & MMLU-Pro & AlpacaEval\\
    \midrule
    \textit{Qwen2.5‑1.5B Results}   \\
    Base & - & 0.002  & 0.090  & 0.000  & 0.000   & 0.297 & 2.10\\
    + GRPO   & MATH  & 0.747  & 0.560  & 0.056  & 0.328   & 0.315 & 4.03\\
    + \methodname  & MATH & 0.711  & 0.530  & 0.099  & 0.296   & 0.310  & 4.28\\
    \midrule
    \textit{Qwen2.5‑3B Results} \\
    Base & - & 0.673  & 0.544  & 0.093  & 0.236   & 0.377 & 3.72\\
    + GRPO & MATH & 0.826  & 0.636  & 0.085  & 0.341   & 0.403 &     6.91\\
    + GRPO-PV  & MATH & 0.820  & 0.636  & 0.086  & 0.299   & 0.398  & 6.17\\
    + \methodname  & MATH     & 0.792  & 0.612  & 0.153   & 0.416   & 0.379  & 7.10\\
    + \methodname-Code  & Codeforces  & 0.743  & 0.572  & 0.153   & 0.411   & 0.386  & 4.16\\
    \bottomrule
  \end{tabular}
   }
   \vspace{-1em}
\end{table}

% \begin{table}[t]
%   \centering
%   \caption{New testing, using 7B system prompt and verl. - indicate original result for base model, to save time I didn't test it another time.}
%   \label{tab:big}
%   \resizebox{0.98\linewidth}{!}{
%   \begin{tabular}{lcccccccc}
%     \toprule
%     Model           & Training Data & GSM8K  & MATH500 & LCB    & CRUX & MMLU-Pro & AlpacaEval\\
%     \midrule
%     \textit{Qwen2.5‑1.5B Results}   \\
%     Base & - & 0.00  & 0.006  & 0.000  & 0.000   & -0.297 & 0.05\\
%     + GRPO   & MATH  & 0.743  & 0.542  & 0.056  & 0.158  & 0.316 & 3.57\\
%     + \methodname  & MATH & 0.705  & 0.528  & 0.059  & 0.300   & 0.314  & 3.88\\
%     \midrule
%     \textit{Qwen2.5‑3B Results} \\
%     Base & - & 0.681  & 0.564  & 0.091  & 0.239   & -0.377 & 4.37\\
%     + GRPO & MATH & 0.842  & 0.658  & 0.169  & 0.408   & 0.403 & 9.86\\
%     + \methodname  & MATH  & 0.806  & 0.616  & 0.154   & 0.219   & 0.386  & 8.43\\
%     \bottomrule
%   \end{tabular}
%    }
% \end{table}


% \begin{table}[t]
%   \centering
%   \caption{New testing}
%   \label{tab:big}
%   \resizebox{0.98\linewidth}{!}{
%   \begin{tabular}{lcccccccc}
%     \toprule
%     Model           & Training Data & GSM8K  & MATH500 & LCB    & CRUX & MMLU-Pro & AlpacaEval\\
%     \midrule
%     + GRPO-new-run0  & MATH   & 0.832  & 0.600  & 0.093   & 0.341   & -  & - \\
%     + GRPO-new-run1  & MATH   & 0.828  & 0.570  & 0.092   & 0.333   & -  & - \\
%     + \methodname-new  & MATH     & 0.803  & 0.636  & 0.118   & 0.386   & -  & - \\
%     + \methodname-new-run0  & MATH     & 0.795  & 0.606  & 0.137   & 0.415   & -  & - \\
%     + \methodname-new-run1  & MATH     & 0.811  & 0.612  & 0.146 & 0.388   & -  & - \\
%     \bottomrule
%   \end{tabular}
%    }
% \end{table}


\section{Results and Analysis}


\begin{figure}[t]
    \centering
    \begin{minipage}[H]{0.495\linewidth}
        \centering
    \includegraphics[width=0.95\linewidth]{fig/completion_length_plots.pdf}
    \vspace{-1.em}
    % \caption{Average response lengths during training rollouts. For Qwen2.5-1.5B, \methodname and GRPO reduce gibberish outputs. For Qwen2.5-3B, \methodname and GRPO increase reasoning length; \methodname yields significantly longer responses. GRPO-PV shows minimal length increase.}
    \caption{\textbf{\methodname encourages longer, more detailed reasoning during training.} This figure shows the average response length during training rollouts on the MATH dataset. For Qwen2.5-1.5B, \methodname and GRPO reduce gibberish outputs. For Qwen2.5-3B, \methodname and GRPO increase reasoning length; \methodname yields significantly longer responses. GRPO-PV shows minimal length increase.}
    \label{fig:length}
    \end{minipage}%
        \vspace{-1.em}
    \hfill
    \begin{minipage}[H]{0.495\linewidth}
    \centering
    \includegraphics[width=0.95\linewidth]{fig/qwen_math_lcb_plots.pdf}
    \vspace{-1.em}
    % \caption{Performance evolution on MATH500 (in-domain) and LiveCodeBench (transfer) for models trained on MATH. MATH500 accuracy increases rapidly at first, preceding gains in code-generation accuracy. LiveCodeBench performance continues to rise even after MATH500 accuracy plateaus.}
    \caption{\textbf{Mastery of in-domain skills facilitates subsequent generalization to new domains.} This figure plots the performance evolution on MATH500 (in-domain, left) and LiveCodeBench (out-of-domain, right) for models trained on the MATH dataset. MATH500 accuracy increases rapidly at first, preceding gains in code-generation accuracy. LiveCodeBench performance continues to rise even after MATH500 accuracy plateaus.}
    \label{fig:split}
    \end{minipage}
    \vspace{-0.5em}
\end{figure}
In this section, we evaluate the effectiveness of \methodname by addressing the following research questions: 
\begin{itemize}[leftmargin=*, nosep]
\item (RQ1) How does the overall performance of \methodname compare to supervised RLVR methods? 
\item (RQ2) How does intrinsic feedback influence the model's qualitative behavior? 
\item (RQ3) How robust is online self-certainty when used as an intrinsic reward signal during training?
\end{itemize}

Table~\ref{tab:big} presents main evaluation results, and Figure~\ref{fig:length} illustrates response length evolution during training. On in-domain MATH and GSM8K datasets, \methodname and GRPO-PV (both golden-answer-free) achieve performance comparable to GRPO (using golden answers). This aligns with TTRL \citep{zuo2025ttrl}, where plurality voting approximated golden answers without significant performance loss. While \methodname performs slightly worse than GRPO overall, on MATH it produces longer responses and demonstrates markedly improved code generation, suggesting enhanced reasoning capabilities.

\subsection{Learning to Follow Instructions}

\methodname significantly enhances instruction-following. Initially, the pretrained Qwen2.5-1.5B struggles with chat-style prompts, scoring $<$10\% on all chat-template tasks (Table~\ref{tab:big}) and generating repetitive, nonsensical output, which inflates average response lengths (Figure~\ref{fig:length}). Fine-tuning with \methodname sharply reduces such gibberish, decreases completion lengths, and enables non-trivial performance across all evaluated benchmarks. Furthermore, on the MATH dataset, \methodname substantially improves the Length Control Win Rate on AlpacaEval for both Qwen2.5-1.5B and Qwen2.5-3B, surpassing GRPO under identical settings. This demonstrates robust gains in instruction adherence.

\subsection{Fostering Structured Reasoning}
\begin{wraptable}{r}{0.5\linewidth}
  \centering
  \vspace{-1em}
  % \caption{Early-stage performance (training step 10) on \textsc{GSM8K} and \textsc{MATH}. \methodname consistently outperforms GRPO.}
  \caption{\textbf{\methodname demonstrates faster initial learning compared to GRPO.} This table shows the in-domain performance on GSM8K and MATH after only 10 training steps. In all cases, \methodname achieves higher accuracy than the GRPO baseline, which uses ground-truth rewards.}
  \label{tab:fast}
  \vspace{-0.5em}
  \resizebox{0.98\linewidth}{!}{
  \begin{tabular}{llcc}
    \toprule
    Model & Method & GSM8K & MATH \\
    \midrule
    \multirow{3}{*}{Qwen2.5-1.5B}
      & Baseline & 0.002 & 0.090 \\
      & GRPO     & 0.081 & 0.296 \\
      & \methodname      & \textbf{0.152} & \textbf{0.368} \\
    \midrule
    \multirow{3}{*}{Qwen2.5-3B}
      & Baseline & 0.673 & 0.544 \\
      & GRPO     & 0.758 & 0.596 \\
      & \methodname      & \textbf{0.811} & \textbf{0.618} \\
    \bottomrule
  \end{tabular}
\  }
  \vspace{-1em}
\end{wraptable}
\textbf{Rapid Initial Learning. }
Self-certainty, a continuous and inherently process-aware reward derived from the model's internal assessment across all tokens, contrasts with binary rewards. This internal signal may encourage LLMs to follow more effective learning trajectories. Given comparable final performance between GRPO and \methodname, we assess early-stage learnability by comparing in-domain accuracy at training step 10. As shown in Table~\ref{tab:fast}, \methodname consistently outperforms GRPO on both GSM8K and MATH benchmarks for Qwen2.5-1.5B and Qwen2.5-3B, highlighting its advantage in rapid initial learning.




\textbf{Cross-Task Generalization. }
Figure~\ref{fig:split} illustrates performance trajectories on MATH500 (in-domain) and LiveCodeBench (transfer task) for models trained on the MATH dataset. For both \methodname and GRPO, accuracy improvements on LiveCodeBench emerge later in training, following initial gains on MATH500. Notably, LiveCodeBench performance continues to improve even after MATH500 accuracy plateaus. This pattern suggests that initial in-domain learning (on MATH) facilitates subsequent generalization to code generation tasks (LiveCodeBench).







\begin{wrapfigure}{r}{0.42\linewidth}
  \vspace{-1em}
  \centering
  \includegraphics[width=0.9\linewidth]{fig/output.pdf}
  % \vspace{em}
  \caption{\methodname quickly demonstrate R1-like reasoning}
  \label{fig:output}
  \vspace{-2em}
\end{wrapfigure}
\textbf{Emergence of Long-Form Reasoning. } While large models like Deepseek-R1 achieve long-form reasoning through extensive RL, \methodname enables smaller models to develop structured reasoning with limited data. On CRUXEval-O (Figure~\ref{fig:output}), models trained with \methodname often exhibit free-form reasoning before summarizing it within the instructed JSON block, despite prompts requiring reasoning directly in JSON. A similar pattern of pre-code natural language reasoning is observed on LiveCodeBench. This emergent pre-reasoning may contribute to \methodname’s strong performance on these benchmarks.

\subsection{Understanding Emergent Long-Form Reasoning}
When LLMs encounter unfamiliar questions, they sample from a distribution of possible answers \citep{kang2024unfamiliar}. Self-certainty reflects the model's internal assessment of its output coherence. By reinforcing high-confidence responses, \methodname encourages more elaborate reasoning, potentially improving the model's comprehension of its own outputs. While not explicitly targeting benchmark accuracy, this enhancement in output quality and structure leads to more reliable answers and better generalization.

We analyze models trained with \methodname on code corpora by examining outputs for ten randomly selected LiveCodeBench questions across different training steps. Figure~\ref{fig:instruction} shows the evolution of output types alongside model accuracy. The results reveal a clear progression: models first learn to generate valid Python code (evidenced by improved accuracy and fewer invalid responses), then develop pre-code reasoning to facilitate self-understanding. Further inspection of generations confirms that models progressively elaborate their reasoning throughout training, supporting our hypothesis that \methodname encourages traces that the model itself can better understand.

To quantify this effect, we classify outputs from successive checkpoints into three categories: invalid code ("No Answer"), valid code without reasoning ("No Reasoning"), and valid code with explicit reasoning ("Reasoning"). Figure~\ref{fig:instruction}(a) illustrates how these proportions evolve during training alongside LiveCodeBench accuracy. The model first reduces invalid outputs and improves code correctness before incorporating pre-code reasoning, reflecting an emergent emphasis on self-explanatory traces. Figure~\ref{fig:instruction}(b) demonstrates how training with \methodname leads to structured reasoning before code generation. Additional evidence appears in Figure~\ref{fig:hist}, where \methodname-trained models assign significantly higher confidence to their generated responses compared to baseline models, as discussed further in Section~\ref{sec:of}.
\begin{figure}[h]
    \centering
    \begin{minipage}[H]{0.58\linewidth}
        \centering
        \includegraphics[width=\linewidth]{fig/ratio_accuracy_plot.pdf}
    \end{minipage}\hfill % \hfill provides horizontal spacing between the minipages
    \begin{minipage}[H]{0.38\linewidth}
        \centering
        \includegraphics[width=\linewidth]{fig/code_prompt.pdf}
    \end{minipage}
    \vspace{-0.8em}
    \caption{(a) Left: Distribution of answer types for ten random LiveCodeBench questions across training steps. Right: Corresponding model accuracy. The model first learns to generate correct code, then adds reasoning to improve understanding. (b) Training with \methodname on code corpora leads to spontaneous reasoning before coding and explanation of outputs.}
    \vspace{-1.7em}
    \label{fig:instruction} 
\end{figure}


\subsection{Online Self-Certainty Prevents Reward Exploitation}\label{sec:of}
Over-optimization against static reward models is a known failure mode in reinforcement learning~\citep{gao2023scaling}. To assess the robustness of self-certainty as a reward, we compare offline self-certainty (rewards from a fixed base model) with online self-certainty (rewards from the evolving policy model), using a reduced batch size of 224 responses per gradient update.

Figure~\ref{fig:of} demonstrates that the offline annotator is susceptible to exploitation. Around the 100th update step, the policy model learns to inflate its self-certainty reward by appending an auxiliary, already-solved problem to its answer for the given question. This exploitation manifests as a sharp increase in response length (dashed line) and a concurrent collapse in validation accuracy (solid line). In contrast, the online annotator, whose reward signal co-evolves with the policy, prevents such reward hacking and maintains stable training.

To further evaluate the quality of self-certainty as a reward signal, we analyze the distribution of self-certainty scores from policies trained with \methodname and GRPO on MATH500 responses (Figure~\ref{fig:hist}). We employ Mann–Whitney U tests to determine if correct responses achieve significantly higher self-certainty scores than incorrect ones. Both GRPO and \methodname models exhibit significantly higher average self-certainty scores, indicating that GRPO also enhances the model's self-assessment capabilities. Notably, policies trained with online self-certainty (i.e., \methodname) show no signs of reward hacking. The \methodname policy yields the lowest $p$-values and largest effect sizes ($r$) in the Mann-Whitney U tests (Figure~\ref{fig:hist}, inset). This indicates it is most effective at discriminating its own correct and incorrect answers using self-certainty, even while assigning higher absolute confidence scores overall. These findings underscore the potential of \methodname for robust training on larger datasets.

\begin{figure}[t]
    \centering
    \vspace{-1em}
    \begin{minipage}[H]{0.46\linewidth}
        \centering
        \includegraphics[width=\linewidth]{fig/metrics.pdf}
        \vspace{-0.5em}
        % \captionof{figure}{Accuracy and response length during training, comparing online and offline self-certainty annotators with \methodname under reduced batch sizes. The offline reward model is exploited early in training (around 100 steps), leading to increased response length and decreased accuracy. The online annotator maintains stable training. Refer to Section~\ref{sec:of} for details.}
        \captionof{figure}{\textbf{Online self-certainty is robust to reward exploitation, unlike offline rewards.} The figure compares the training stability of \methodname with an online self-certainty annotator (updated with the policy) versus an offline one (fixed base model). The policy rapidly exploits the static offline annotator, causing a spike in response length and a drop in accuracy near step 100. In contrast, the evolving online reward avoids exploitation and enables stable training (Sec.~\ref{sec:of}).}
        \label{fig:of}
    \end{minipage}%
    \hfill
    \begin{minipage}[H]{0.53\linewidth}
        \centering
        \includegraphics[width=\linewidth]{fig/confidence_histograms.pdf}
        \vspace{-2em}
        % \captionof{figure}{Distribution of self-certainty on \textsc{MATH500} responses, for policies trained with GRPO and \methodname. Histograms are split by response correctness. The inset shows Mann–Whitney U test statistics ($p$-value and effect size $r$) comparing self-certainty of correct versus incorrect responses. The policy trained with \methodname demonstrates the best separation.}
        \captionof{figure}{\textbf{Training with \methodname improves the model's ability to distinguish its own correct and incorrect answers.} Distributions of self-certainty on \textsc{MATH500} are shown for policies trained with GRPO and \methodname. Histograms are split by response correctness. The inset shows Mann–Whitney U test statistics ($p$-value and effect size $r$) comparing self-certainty of correct versus incorrect responses. The policy trained with \methodname demonstrates the best separation.}
        \label{fig:hist}
    \end{minipage}
    \vspace{-1.3em}
\end{figure}

\subsection{Ablation Studies}
To comprehensively validate \methodname's design and robustness, we conducted extensive ablation studies, with full details provided in Appendix \ref{sec:app_exp} due to page limitations. Key findings are: (1) KL term: Varying the KL penalty (Sec. \ref{sec:app_kl}) shows a stability–performance trade-off; moderate values yield the best accuracy.
(2) Scaling: \methodname scales to larger backbones (Qwen2.5-7B/14B, Qwen3-14B; Sec. \ref{sec:app_large_models}), delivering consistent gains in reasoning and generalization.
(3) Architecture: On Llama-3.2 and OLMo-2 (Sec. \ref{sec:llama}), \methodname remains effective, indicating robustness across model families and sizes.
(4) Reward design: Compared to entropy minimization \citep{agarwal2025unreasonable} and random rewards \citep{shao2025spurious}, \methodname yields stable improvements, while the alternatives trigger catastrophic collapse (Sec. \ref{sec:add_reward}).
(5) Optimization strategy: Directly optimizing self-certainty as a loss function leads to reward hacking and performance collapse; our advantage-weighted policy-gradient formulation avoids this and trains reliably (Sec. \ref{sec:app_loss}).

\section{Discussion and Future Research}

\textbf{Scalability and Generalization. }
Our experiments, constrained by computational resources, utilize relatively compact models trained on relatively small, unsupervised corpora. We aim to demonstrate the potential of a model's self-certainty as a reward signal for policy optimization. The results show that this signal consistently promotes more coherent, well-justified, and interpretable explanations, indicating a path towards more autonomous learning. Future work could explore these benefits in larger foundation models (with hundreds of billions of parameters) and on more diverse, real-world datasets. Given that purely offline training with \methodname led to performance degradation over time, scaling up will likely require periodic online updates to self-certainty estimates or hybrid offline-online schedules to maintain calibration.

% \textbf{Applicability to Other Policy Gradient Methods. }
% \methodname is a framework that leverages a model's self-certainty as an intrinsic reward signal for fine-tuning LLMs. It can be instantiated with various policy gradient algorithms. Due to computational constraints, and informed by the success of models like DeepSeek-R1 and the widespread adoption of GRPO, we employ GRPO for self-certainty optimization. The efficacy of self-certainty signals with other algorithms, such as REINFORCE or PPO, warrants further investigation.

\textbf{Theoretical Analysis of RLIF. } While we have empirically demonstrated the superior performance of using self-certainty as a reward for RLIF, the underlying theoretical mechanisms warrant further investigation. \citet{huang2025self} analyze LLM self-improvement as a ``sharpening'' mechanism, proposing a statistical framework to evaluate algorithm efficiency via sample complexity. Similarly, \citet{yue2025does} question whether RLVR functions primarily by sharpening the base model's existing distribution. However, determining the theoretically optimal reward signal for RLIF and establishing the fundamental reasoning boundaries of LLMs remain open problems. These challenges highlight the need for future theoretical research to complement empirical findings.

\textbf{Combining Reward Signals. }
To enable a direct comparison between self-certainty and golden-answer rewards, this paper focuses exclusively on a single reward signal. However, these signals are not mutually exclusive. Future work could explore combining them, for instance, by summation or by alternating based on the availability of golden answers. Furthermore, other reward signals, such as formatting rewards \citep{guo2025deepseek}, could be additively combined to enhance performance. Integrating RLIF with methods like RLHF and RLVR may further advance LLM capabilities across various dimensions.


\section{Conclusion}
% This paper introduces \methodname, a reinforcement learning framework using intrinsic self-certainty as a reward, eliminating the need for external supervision. \methodname matches GRPO's performance on mathematical reasoning tasks without gold solutions and demonstrates significantly better generalization to downstream tasks like code generation. It also enhances instruction-following and fosters structured reasoning, with online self-certainty proving robust against reward exploitation. These results highlight the potential of intrinsic signals for effective, scalable, and generalizable LLM fine-tuning. 
This paper introduces \methodname, an instantiation of Reinforcement Learning from Internal Feedback (RLIF) that uses a model’s intrinsic self-certainty as its sole reward signal, eliminating the need for external supervision or gold-standard solutions. Our experiments show that \methodname matches the performance of supervised RLVR methods like GRPO on mathematical reasoning and achieves competitive, sometimes better generalization to out-of-domain tasks such as code generation and instruction following. It also promotes structured reasoning and leverages online self-certainty to guard against reward exploitation. 


These findings highlight the transformative potential of RLIF, signaling a meaningful step toward AI systems that improve through introspection and unlock rich latent capabilities. Looking forward, this paradigm opens the door to AI agents capable of autonomous skill acquisition in novel domains and scalable self-improvement—even as they approach or surpass the limits of human oversight. Future directions include integrating RLIF with external reward methods like RLHF or RLVR to tackle increasingly complex real-world challenges, and advancing the development of more robust, generalizable, and truly autonomous learning systems.

\section*{Ethics Statement}
Our research is based on publicly available datasets and open-source language models, mitigating concerns related to private data or human subjects. The goal of our work is to enhance the reasoning capabilities of language models through self-supervision, which we believe is a positive step toward more transparent and robust AI systems. We have made our code publicly available to ensure transparency and allow for full scrutiny of our methods and findings. We do not foresee any direct negative societal impacts or ethical concerns arising from this work.

\section*{Reproducibility Statement}
To ensure the reproducibility of our results, we provide all source code and training configurations in \url{https://github.com/sunblaze-ucb/Intuitor}. The Experimental Setup section and Appendix \ref{sec:app_exp} detail all hyperparameters, software versions (including the Open-R1 framework), and evaluation setups. Furthermore, Appendix \ref{app:prompt-completions} includes the exact prompts used during training and evaluation. These resources should allow for the complete replication of our experiments and validation of our findings.

\bibliography{ref}
\bibliographystyle{iclr2026_conference}

\appendix
\newpage

\input{appendix}


\end{document}
